\documentclass[12pt, a4paper]{article}

% ---- Packages ----
\usepackage[utf8]{inputenc}
\usepackage[T1]{fontenc}
\usepackage{amsmath, amssymb}
\usepackage{graphicx}
\usepackage{booktabs}
\usepackage{hyperref}
\usepackage{geometry}
\usepackage{caption}
\usepackage{subcaption}
\usepackage{float}
\usepackage{enumitem}
\usepackage{xcolor}
\usepackage{listings}
\usepackage{fancyhdr}

\geometry{margin=2.5cm}

% ---- Code listing style ----
\lstset{
    language=Matlab,
    basicstyle=\ttfamily\small,
    keywordstyle=\color{blue}\bfseries,
    commentstyle=\color{green!60!black},
    stringstyle=\color{red!70!black},
    numbers=left,
    numberstyle=\tiny\color{gray},
    numbersep=8pt,
    frame=single,
    breaklines=true,
    backgroundcolor=\color{gray!5},
    tabsize=4,
    showstringspaces=false
}

% ---- Header / Footer ----
\pagestyle{fancy}
\fancyhf{}
\rhead{CMPE 362 -- HW3}
\lhead{Ali Ayhan Günder - 2021400219}
\cfoot{\thepage}

% ---- Title ----
\title{
    \vspace{-1cm}
    \textbf{CMPE 362 -- Homework 3} \\[0.3cm]
    \Large Aliasing in Audio, Image, and Video
}
\author{
    Ali Ayhan Günder \\
    Student ID: 2021400219 \\[0.2cm]
    \textit{CMPE 362 -- Signal Processing, Spring 2026}
}
\date{\today}

\begin{document}

\maketitle
\thispagestyle{fancy}

\begin{abstract}
This report investigates aliasing across three signal domains---audio, image, and video---demonstrating that the Nyquist--Shannon Sampling Theorem governs all three, despite their different perceptual manifestations. Through MATLAB experiments, we show how undersampling creates structured phantom signals: a false 200\,Hz tone from a 1800\,Hz source, Moir\'{e} patterns from a 2D chirp function, and the wagon wheel illusion from temporal undersampling of rotational motion.
\end{abstract}

\tableofcontents
\newpage

%==========================================================================
\section{Introduction}
%==========================================================================

Aliasing is a fundamental artifact that arises when a continuous signal is sampled at a rate insufficient to capture its highest frequency content. When the sampling rate violates the Nyquist criterion, high-frequency components do not simply vanish---they \emph{masquerade} as lower-frequency components, producing deterministic phantom signals that are indistinguishable from genuine low-frequency content.

This report investigates aliasing across three signal domains:
\begin{enumerate}[label=(\alph*)]
    \item \textbf{Audio:} A 1800\,Hz cosine tone sampled below and above the Nyquist rate.
    \item \textbf{Image:} A 2D chirp function $\cos\!\big((x+2y)^2\big)$ rendered at decreasing resolutions.
    \item \textbf{Video:} The wagon wheel effect as temporal aliasing of rotational motion.
\end{enumerate}

We demonstrate that all three phenomena are governed by the same mathematical principle.

%==========================================================================
\section{Theoretical Background}
%==========================================================================

\subsection{The Nyquist--Shannon Sampling Theorem}

A band-limited continuous signal with maximum frequency $f_{\max}$ can be perfectly reconstructed from its samples if and only if:
\begin{equation}
    f_s > 2 \cdot f_{\max}
    \label{eq:nyquist}
\end{equation}
where $f_s$ is the sampling rate. The critical threshold $f_N = f_s / 2$ is called the \textbf{Nyquist frequency}.

\subsection{Frequency Folding}

When a signal of frequency $f$ is sampled at rate $f_s < 2f$, the sampled signal appears as a different, lower-frequency signal. The apparent (alias) frequency is:
\begin{equation}
    f_{\text{alias}} = \left| f - k \cdot f_s \right|, \quad k = \text{round}(f / f_s)
    \label{eq:alias-general}
\end{equation}

For the common case where $f_s/2 < f < f_s$:
\begin{equation}
    f_{\text{alias}} = f_s - f
    \label{eq:alias-simple}
\end{equation}

This ``frequency folding'' around $f_N$ is the core mechanism behind all aliasing phenomena examined in this report.

%==========================================================================
\section{Part A -- Audio Aliasing}
%==========================================================================

\subsection{Experiment Setup}

We generate a pure cosine tone at $f = 1800$\,Hz and sample it at two different rates (2000\,Hz and 4000\,Hz). Additionally, a reference tone at 200\,Hz sampled at 2000\,Hz is generated for comparison.

\begin{table}[H]
\centering
\caption{Audio aliasing experiment configurations.}
\label{tab:audio-config}
\begin{tabular}{@{}clcccc@{}}
\toprule
\textbf{Config} & \textbf{Signal} & $f_s$ \textbf{(Hz)} & $f_N$ \textbf{(Hz)} & \textbf{Nyquist OK?} & $f_{\text{alias}}$ \textbf{(Hz)} \\
\midrule
A & $\cos(2\pi \cdot 1800\,t)$ & 4000 & 2000 & Yes & 1800 \\
B & $\cos(2\pi \cdot 1800\,t)$ & 2000 & 1000 & No  & 200  \\
Ref & $\cos(2\pi \cdot 200\,t)$ & 2000 & 1000 & Yes & 200  \\
\bottomrule
\end{tabular}
\end{table}

\subsection{MATLAB Implementation}

\begin{lstlisting}[caption={Audio signal generation and spectrogram computation.}]
f = 1800; dur = 2;

% Config A: 1800 Hz @ fs=4000 -> Properly sampled
fs2 = 4000; t2 = 0:1/fs2:dur-1/fs2;
y2 = cos(2*pi*f*t2);

% Config B: 1800 Hz @ fs=2000 -> Aliased to 200 Hz
fs1 = 2000; t1 = 0:1/fs1:dur-1/fs1;
y1 = cos(2*pi*f*t1);

% Reference: 200 Hz @ fs=2000
y_low = cos(2*pi*200*t1);
\end{lstlisting}

Spectrograms are computed using a custom Short-Time Fourier Transform (STFT) implementation with a 256-sample Hamming window and 50\% overlap ($n_{\text{overlap}} = 128$).

\subsection{Spectrogram Results}

\begin{figure}
    \centering
    \includegraphics[width=1.0\linewidth]{\detokenize{Audio Aliasing Spectrograms.png}}
    \caption{Audio spectrograms. \textbf{Top:} 200\,Hz at $f_s=2000$\,Hz. \textbf{Middle:} 1800\,Hz at $f_s=2000$\,Hz (aliased to 200\,Hz). \textbf{Bottom:} 1800\,Hz at $f_s=4000$\,Hz.}
    \label{fig:audio-spectrograms}
\end{figure}

\subsection{Analysis}

\paragraph{Config A ($f_s = 4000$\,Hz):}
The Nyquist condition is satisfied ($1800 < 2000$). The spectrogram shows a clean horizontal band at 1800\,Hz, faithfully representing the original tone.

\paragraph{Config B ($f_s = 2000$\,Hz):}
The alias frequency is:
\begin{equation}
    f_{\text{alias}} = f_s - f = 2000 - 1800 = 200 \text{\,Hz}
\end{equation}
The sampled waveform is \textbf{mathematically identical} to a genuine 200\,Hz cosine. The spectrogram shows a band at 200\,Hz---indistinguishable from the reference. When played back, both sound like the same low-pitched hum.

\paragraph{Perceptual consequence:} A listener cannot distinguish whether the source was a 1800\,Hz tone sampled too slowly or a genuine 200\,Hz tone. The original high-frequency information is \textbf{irrecoverably lost}.

%==========================================================================
\section{Part B -- Image Aliasing}
%==========================================================================

\subsection{Experiment Setup}

We render the 2D function:
\begin{equation}
    I(x, y) = \cos\!\Big((x + 2y)^2\Big), \quad x \in (-10, 10), \; y \in (-10, 10)
    \label{eq:chirp}
\end{equation}

This is a \textbf{2D chirp signal}: the argument $(x + 2y)^2$ means the instantaneous spatial frequency increases quadratically as we move away from the line $x + 2y = 0$. Near that line, the pattern has low frequency (wide bands); far from it, the frequency is extremely high (narrow bands).

We sample this function at four resolutions: $2000 \times 2000$, $500 \times 500$, $150 \times 150$, and $50 \times 50$.

\subsection{MATLAB Implementation}

\begin{lstlisting}[caption={Image generation at multiple resolutions.}]
res_levels = [2000, 500, 150, 50];
images = cell(1, length(res_levels));

for i = 1:length(res_levels)
    res = res_levels(i);
    [x, y] = meshgrid(linspace(-10, 10, res), ...
                       linspace(-10, 10, res));
    images{i} = cos((x + 2*y).^2);
end
\end{lstlisting}

\subsection{Results}

\begin{figure}
    \centering
    \includegraphics[width=1.0\linewidth]{\detokenize{Image Aliasing.png}}
    \caption{Image samples of $I(x,y)=\cos((x+2y)^2)$ at four resolutions ($2000\times2000$, $500\times500$, $150\times150$, $50\times50$). As resolution decreases, Moir\'{e} patterns (phantom low-frequency structures) appear due to spatial aliasing.}
    \label{fig:image-aliasing}
\end{figure}

\begin{table}[H]
\centering
\caption{Observations at each resolution level.}
\label{tab:image-obs}
\begin{tabular}{@{}ll@{}}
\toprule
\textbf{Resolution} & \textbf{Observation} \\
\midrule
$2000 \times 2000$ & Clean chirp---hyperbolic fringes centered on $x + 2y = 0$, thinning toward corners. \\
$500 \times 500$ & Center sharp; faint false patterns appear at high-frequency edges. \\
$150 \times 150$ & \textbf{Moir\'{e} patterns emerge dramatically.} Phantom large-scale rings visible. \\
$50 \times 50$ & Entirely dominated by aliased macro-patterns. True structure obscured. \\
\bottomrule
\end{tabular}
\end{table}

\subsection{Analysis}

The spatial sampling rate is $f_s = \text{resolution} / 20$ samples per unit length (since $x$ spans 20 units). The local spatial frequency of $\cos((x+2y)^2)$ along the gradient direction $\nabla(x+2y)^2 = 2(x+2y)(1, 2)$ is:
\begin{equation}
    f_{\text{local}} = \frac{1}{2\pi} \cdot 2|x + 2y| \cdot \sqrt{1^2 + 2^2} = \frac{2\sqrt{5}}{\pi}|x + 2y|
\end{equation}

This frequency increases without bound as $|x + 2y|$ grows. Therefore, no matter how high the resolution, there will always be a region near the corners where the spatial Nyquist limit is violated. At lower resolutions, the Nyquist boundary shrinks toward the center, and high-frequency content everywhere else \textbf{folds back} into low-frequency Moir\'{e} patterns.

\textbf{Key insight:} The Moir\'{e} patterns are not noise or blur---they are \textbf{structured false signals} created by frequency folding, exactly analogous to the 200\,Hz phantom tone in the audio experiment.

%==========================================================================
\section{Part C -- Video Aliasing (Wagon Wheel Effect)}
%==========================================================================

\subsection{The Phenomenon}

In film or video, a spinning wheel sometimes appears to:
\begin{itemize}
    \item Rotate \textbf{backwards} (wrong direction),
    \item Appear \textbf{stationary} (frozen), or
    \item Rotate at the \textbf{wrong speed}.
\end{itemize}
This is the temporal equivalent of the audio and image aliasing demonstrated above.

\subsection{Mapping to Sampling Theory}

\begin{table}[H]
\centering
\caption{Aliasing concepts mapped across domains.}
\label{tab:domain-map}
\begin{tabular}{@{}lll@{}}
\toprule
\textbf{Concept} & \textbf{Audio Domain} & \textbf{Video Domain} \\
\midrule
Continuous signal   & Sound wave                  & Continuous wheel rotation \\
Signal frequency    & Tone frequency $f$ (Hz)     & Spoke angular rate (rev/s) \\
Sampling rate       & $f_s$ (samples/s)           & Frame rate (FPS) \\
Nyquist frequency   & $f_s / 2$                   & FPS / 2 \\
Aliased output      & False lower-pitched tone    & False rotation direction/speed \\
\bottomrule
\end{tabular}
\end{table}

\subsection{Mathematical Description}

Consider a wheel with $N$ evenly spaced spokes, rotating at $R$ revolutions per second, filmed at $F$ frames per second. The angular displacement of the nearest spoke per frame is:
\begin{equation}
    \Delta\theta = \frac{R}{F} \pmod{1/N}
\end{equation}

When $\Delta\theta > \frac{1}{2N}$ (i.e., the spoke moves more than half a spoke-spacing between frames), temporal aliasing occurs:

\begin{itemize}
    \item If a spoke advances by \emph{almost} one full spoke-spacing between frames, it appears to have moved \textbf{backwards by a tiny amount}---the wheel appears to reverse.
    \item If it advances \emph{exactly} one spoke-spacing, the wheel appears \textbf{stationary}.
    \item If it advances \emph{slightly more} than one spacing, it appears to rotate \textbf{slowly forward}.
\end{itemize}

\subsection{The Negative Frequency Analogy}

In the audio experiment, 1800\,Hz sampled at 2000\,Hz aliases to 200\,Hz. In the discrete-time Fourier domain, the alias sits at $-200$\,Hz on one side of the spectrum---a \textbf{negative frequency}. For real-valued audio signals, positive and negative frequencies are conjugate-symmetric:
\begin{equation}
    X(-f) = X^*(f)
\end{equation}
so $-200$\,Hz sounds identical to $+200$\,Hz. Humans perceive only the magnitude, not the sign.

In video, however, \textbf{direction is visible}. The negative frequency manifests as \textbf{reverse rotation}. This is why the wagon wheel effect is such a powerful demonstration of aliasing---it lets us \emph{see} the sign of the aliased frequency that is hidden in audio.

%==========================================================================
\section{Cross-Domain Comparison}
%==========================================================================

\begin{table}[H]
\centering
\caption{Unified comparison of aliasing across all three domains.}
\label{tab:cross-domain}
\renewcommand{\arraystretch}{1.3}
\begin{tabular}{@{}p{2.8cm}p{3.5cm}p{3.5cm}p{3.5cm}@{}}
\toprule
\textbf{Aspect} & \textbf{Audio} & \textbf{Image} & \textbf{Video} \\
\midrule
Signal & $\cos(2\pi \cdot 1800\,t)$ & $\cos((x{+}2y)^2)$ & Spoke angular position \\
Sample rate & 2000\,Hz (temporal) & 150\,px/row (spatial) & 24 FPS (temporal) \\
Aliasing artifact & Phantom 200\,Hz tone & Moir\'{e} fringes & Backward rotation \\
Perception & Hear wrong pitch & See false patterns & See wrong direction \\
Root cause & $f > f_s/2$ & $f_{\text{spatial}} > f_{N,\text{spatial}}$ & $f_{\text{spoke}} > \text{FPS}/2$ \\
Reversible? & No & No & No \\
\bottomrule
\end{tabular}
\end{table}

\textbf{Unified principle:} In every domain, when the signal's frequency exceeds half the sampling rate, the sampled representation folds the true frequency into a false lower one. The artifact is not random noise---it is a deterministic, structured phantom signal.

%==========================================================================
\section{Conclusion}
%==========================================================================

This homework demonstrated aliasing across three signal modalities:

\begin{enumerate}
    \item \textbf{Audio:} A 1800\,Hz tone sampled at 2000\,Hz becomes indistinguishable from 200\,Hz, confirming the alias formula $f_{\text{alias}} = f_s - f$.
    \item \textbf{Image:} The chirp function $\cos((x+2y)^2)$ produces striking Moir\'{e} patterns when undersampled, as high spatial frequencies fold into phantom low-frequency structures.
    \item \textbf{Video:} The wagon wheel effect visually reveals what audio hides---the \emph{sign} of the aliased frequency---because reverse rotation corresponds to a negative frequency component.
\end{enumerate}

All three phenomena are governed by the same Nyquist--Shannon theorem. Aliasing is not domain-specific; it is a universal consequence of insufficient sampling.

\end{document}
